\cotefile{156-8}
\scanpage{66}\mspage{57}

Donc \unsure{finalement} \unsure{je} retombe sur
\[
 f_M=f\mathbin{:}M\longrightarrow M'
\]
comme \unsure{pages|page} 27, 28. On a vu \inserted{alors} que pour \deleted{tout} $\widetilde F\in\unsure{\mathcal{F}\mathrm{ig}|\mathrm{Fig}}(M)$ (figure associée au magasin), \deleted{\unsure{posant}} $F'=f(F)$, on a
\[
 f\mid\widetilde F\mathbin{:}\quad \widetilde F\xrightarrow{\sim}\widetilde F'
 \qquad\begin{gathered}\text{\unsure{isom.|hom.|Isom.} d'ens.}\\[-2pt]\text{ordonnés}\end{gathered}
\]
ce qui donne aussitôt (\unsure{$**$|$xx$}). On dispose \unsure{donc|de} \deleted{\unsure{ce}} de tous les corollaires \unsure{des|de la} p. 28, 29.

Les propriétés de $f$ sont voisines de celles d'un plongement, mais \deleted{en général} $f$ n'est \unsure{pourtant} pas \unsure{néc.|nec.} injectif. Exemple: \unsure{si} $M$ est quasi-ensembliste et non ensembliste, l'hom. canonique
\[
 M\longrightarrow\mathrm{Figures}\,(\unsure{L|\mathcal L})\quad .
\]
Si c'est injectif, est-ce qu'on peut identifier $M{\unsure{{}'|\vphantom{'}}}$ à un \uline{sous-\deleted{\unsure{ma}}\unsure{magasin|catégorie}} de $M'$, i.e. l'image est-elle un sous-magasin \unsure{?}, \inserted{et $M$ y est-il isomorphe} \unsure{Pour|Par} ce qui précède, \deleted{\unsure{elle}} \inserted{l'image} est fermée pour $\leq$. \deleted{\unsure{Montrons que}} \inserted{Donc} si
\[
 f(X)\mathrel{\overset{\circ}{\ll}}Y'\leq f(Y)\qquad\text{dans }M'
\]
\unsure{est-que|est-ce que} $Y'$ est de la forme $f(Y_1)$ ? On voudrait de plus, bien \unsure{sur|sûr}
\[
 X\mathrel{\overset{\unsure{\circ|\cdot}}{\ll}}Y_1\leq Y
\]
en résumé, on voudrait que
% archive sheet number: 66
\scanpage{67}\mspage{58}
\begingroup\csname tagsleft@true\endcsname

\[
 f(X)\ll f(Y)\implies X\ll Y \tag*{(\unsure{$*xx$|$***$})}
\]
ce qui n'est pas automatique je présume, \unsure{$=$|même} si $f$ injectif. Si c'est vrai, alors \unsure{prenant} $Y_1$ avec $X\mathrel{\overset{\circ}{\ll}}Y_1\leq Y$, on aura $f(X)\mathrel{\overset{\circ}{\ll}}f(Y_1)\leq f(Y)$ donc $Y'=f(Y_1)$, OK. De plus, si $f(X)\mathrel{\overset{\circ}{\ll}}f(Y)$, on aura $X\ll Y$, et \unsure{.} \deleted{$f(Y_1)=$} on trouve $f(Y_1)=f(Y)$, donc $Y_1=Y$, donc $X\mathrel{\overset{\circ}{\ll}}Y$. Donc (\unsure{$***$|$xxx$}) impliquerait que $\mathrel{\overset{\circ}{\ll}}$ est induit par $\ll$. Qu'en est-il de $\leq$, $|\unsure{\circ|0}|$ ? Il semble qu'il faille le poser en sus\unsure{.... .}

\figure{\begin{tikzpicture}
\draw[line width=.55pt] (0,1.04) -- (.75,0);
\end{tikzpicture}}

J'ai envie maintenant de déterminer les homomorphismes de magasins
\[
 f\mathbin{:}M\longrightarrow\unsure{\text{Figél}|\text{Figéls}|\text{Figélem}}(P)\qquad ,
\]
où $P$ est un ens. Si on a un tel hom, \unsure{on} \unsure{en|on} déduit une application
\[
 \begin{aligned}
 X\longmapsto{} & |X| &&\mathbin{:}\quad M\longrightarrow\unsure{\mathcal P|P}(P)\quad ,\\[-2pt]
 & \Vert\,\unsure{\mathrm{def}|\text{déf}}\\[-2pt]
 & |f(X)|
 \end{aligned}
\]
% slanted
\margin{On \unsure{utilise} ici que $\widetilde X\xrightarrow{\sim}\widetilde{f(\unsure{F|X})}$ via $f$}
et on retrouve $f$ via cette application, car on aura
\[
 f(X)=\{\,\deleted{\unsure{X|Y}}\,|Y|\mid Y\in\widetilde X\,\}\quad . \tag*{(\unsure{$*$|$x$})}
\]
\endgroup
% archive sheet number: 67
\scanpage{68}\mspage{59}

La question est maintenant, quelles sont les conditions sur une application \deleted{\unsure{$f$}}
\[
 X\mathrel{\unsure{\mapstochar\relbar|\mapsto}}|X|\mathbin{:}\quad M\longrightarrow\unsure{\mathcal P|P}(P),
\]
pour \deleted{\unsure{défi}} pouvoir être définie par un morphisme d'ateliers ? Nous voilà revenus à la situation p. 33 ! En termes de $X\mathrel{\unsure{\text{---}|\to}}|X|$, on définira, pour \deleted{$X\in M$} toute \deleted{$\unsure{\partial X}$} \inserted{partie} $\deleted{\unsure{A}}\,\Phi\subset M$
\[
 |\Phi|=\bigcup_{X\in\Phi}|X|
\]
\unsure{Ce|ce} qui donne un sens à
\[
 |\partial X|\unsure{\subset|<}|X|,\quad |\partial X|=\bigcup_{Y<X}|Y|
\]
\[
 \bigl(\text{\unsure{où|ou}}\qquad\partial X=\widetilde X-\{X\}=\{Y\in M\mid Y<X\}\bigr),
\]
et \unsure{.}
\[
 |X|^\circ=|X|\mathbin{\backslash}|\partial X|.
\]

\deleted{\unsure{On}} Il faudra avoir en tous cas
\begingroup\csname tagsleft@true\endcsname
\[
 X\ll Y\implies |X|\leq |Y| \tag*{(i)}
\]
\[
 X\,|\unsure{\circ|0}|\,Y\implies |X|^\circ\cap\deleted{\unsure{-}}|Y|^\circ=\varnothing \tag*{(ii)}
\]
\[
 \forall X\in M,\qquad |X|^\circ\ne\varnothing \tag*{(iii)}
\]
\[
 \forall \unsure{X|x}\in M,\qquad |X|=\bigcup_{Y\unsure{\leq|\in} X}|Y|^\circ \tag*{(iv)}
\]
\endgroup
Montrons que ces conditions permettent déjà de définir, pour \unsure{tt|$\forall$} $X\in M$, un $f(X)\in\text{Fig él}\,(\unsure{.}P)$ par (\unsure{$x$|$*$}) page précédente. Il faut donc voir que la famille
% archive sheet number: 68
\scanpage{69}\mspage{60}

des $|Y|$, $Y\in\widetilde X$, est une figure élémentaire (de support $|X|$, bien sûr), car par (i) les $|Y|\subset |X|$.)\unsure{,|.} Cela résulte en effet du critère (\unsure{ivi|iVi|iii}) de la prop. p. 37. (\uline{NB} on n'a utilisé (i) que \unsure{pour} le cas $X\leq Y$.) Il est clair que cette application $f$ est croissante pour $\leq$,
% slanted
\margin{condition d) p. 28}
\deleted{et que} \inserted{et qu'elle \unsure{induit}, $\forall X\in M$, $\widetilde X\xrightarrow{\sim}\widetilde{f(X)}$. \unsure{Ainsi,|Aussi,}} l'on a bien
\[
 |X|^\circ=f(X)^\circ
\]
de sorte que (ii) implique également
\[
 X\,|\unsure{\circ|0}|\,Y\implies f(X)\,|\circ|\,f(Y).
\]
Il reste à vérifier que
\[
 X\mathrel{\overset{\circ}{\ll}}Y\implies f(X)\mathrel{\overset{\circ}{\ll}}f(Y),
\]
\deleted{Mais} va-t-on y arriver ? On aura une application \uline{croissante}
\[
 \unsure{\phi|\varphi}\mathbin{:}\widetilde X\longrightarrow\widetilde Y,\qquad \varphi(X)=Y
\]
\deleted{telle que} caractérisée par
\[
 X'\mathrel{\overset{\circ}{\ll}}\varphi(Y')
\]
et \deleted{\unsure{tout}} revient à ceci:
% slanted
\margin{cf \unsure{CF|CE}\,$\overline{\underline{\mathrm V}}$ p. \unsure{19}, 20, \unsure{Critère|critère} \unsure{b|3}}
\[
 X\mathrel{\overset{\circ}{\ll}}Y\implies |X|^\circ\subset |Y|^\circ
\]
\deleted{\unsure{On} sait déjà \unsure{q}} Par (i) on aura
\[
 |X|\subset |Y|\qquad\text{d'où}\quad |X|^\circ\subset |Y|,
\]
% archive sheet number: 69
\scanpage{70}\mspage{61}

il reste à voir que
\[
 |X|^\circ\cap|\partial Y|=\varnothing
\]
ce qui revient à
\[
 |X|^\circ\cap|Z|^{\deleted{\unsure{\circ}}}=\varnothing\qquad \forall Z\deleted{\unsure{\in}}\!<Y
\]
\deleted{\unsure{et si on}} \deleted{\unsure{(p}} et comme
\[
 |Z|=\bigcup_{Z'\leq Z}|Z'|^\circ
\]
par (iv), il revient au \unsure{$=$|même} de voir
\[
 |X|^\circ\cap|Z|^\circ=\varnothing\qquad\forall Z<Y
\]
ce qui va résulter de
\[
 X\,|\unsure{\circ|0}|\,Z\qquad\text{(via (ii))}
\]
or si $Z<Y$ on a $Z\,|\circ|\,Y$ ($\unsure{M_\theta\,4|M_\theta^{4}|\mathrm{Mag}\,4}$) et comme $X\mathrel{\overset{\circ}{\ll}}Z$ on en déduit $X\,|\circ|\,Z$ ($M_\theta\unsure{\,3|{}^{3}}$), on gagne !

Donc

\uline{Théorème} Les hom. de magasins$^{*}$
% slanted
\margin{cf p. 27 pour la \unsure{notion|notation}}
\[
 M\longrightarrow\mathrm{Figures}\,(P)\ \unsure{*|,}
\]
sont en corr. \unsure{1-1|1=1} avec les applications
\[
 \begin{aligned}
 M&\longrightarrow \unsure{\mathcal P|P}(P)\\
 X&\mathrel{\unsure{\text{---}|\to}}|X|
 \end{aligned}
\]
satisfaisant \unsure{:|à} les \deleted{applications} \inserted{conditions} (i)\,--\,(iv) de la page \deleted{\unsure{6}}59.

Une fonction ayant ces propriétés est
% archive sheet number: 70
\scanpage{71}\mspage{62}
% archive sheet number: 71

appelée \uline{fonction support externe} (à valeurs dans $\unsure{\mathcal P|P}(P)$). Elle est dite \uline{normalisée} si a) $P=|\deleted{\unsure{X}}M|$ \unsure{i.e|i.e.} les $|X|$ recouvrent $P$ et b) les $|X|$ séparent les points de $P$ \deleted{\unsure{(et alors}} \inserted{($\Longleftrightarrow$} les $|X|^\circ$ \deleted{aussi} séparent les points). Elle est dite \uline{fidèle} si les \deleted{\unsure{conditi}} deux conditions (ii) (iii) \unsure{se renforcent} en la condition
\begingroup\csname tagsleft@true\endcsname
\[
 \forall X,Y\in M,\qquad X\,|\unsure{\circ|0}|\,Y
 \Longleftrightarrow |X|^\circ\cap|Y|^\circ=\varnothing \tag*{(\unsure{Supp fid})}
\]
\endgroup
% slanted
\margin{On dit \deleted{aussi} qu'elle constitue une \uline{``réalisation} \uline{\unsure{``}ensembliste''} si \uline{de plus}, les relations $\mathrel{\overset{\circ}{\ll}}$, $\leq$ sont ``induites'' \deleted{\unsure{i.e.}} \deleted{\unsure{$\Leftrightarrow$}} i.e.
\[
 \begin{aligned}
 \unsure{\varphi|\phi}(X)\mathrel{\overset{\circ}{\ll}}\varphi(Y)&\Rightarrow X\mathrel{\overset{\circ}{\ll}}Y\\
 \varphi(X)\leq\varphi(Y)&\mathrel{\unsure{\nRightarrow|\Rightarrow}}X\leq Y
 \end{aligned}
\]}

\uline{NB} On peut tjrs, en remplaçant \deleted{$X$} \inserted{$P$} par un quotient de $|M|$, se ramener au cas d'une fonction support normalisée.

\uline{Scholie} Dans tous les exemples de \unsure{magasins} que j'ai effleurés \unsure{jusqu'ici}, provenant de contextes divers \unsure{ensemblistes|ensembliste}, de topologie générale ou de topologie modérée, il existe toujours$^{*}$
% slanted
\margin{2.7. C'est inexact -- dans le cas des \unsure{pseudosegments|pseudo-segments} topologiques (ou ``\unsure{cordes|ordres}'' ou ``\unsure{ficelles|feuilles}'') etc, la fonction support \unsure{(naturelle)}, n'est en général \uuline{pas} fidèle, contrairement à ce qui \unsure{m'avait|m'aurait} semblé d'abord -- il est vrai que \unsure{$\varphi$} $\varphi(X)\mathrel{\overset{\circ}{\ll}}\varphi(Y)\Rightarrow X\ll Y$ et \unsure{itou|idem} pour $\ll$, $\leq$,}
des fonctions support \uline{fidèles} sur des ens. convenables. Donc les \deleted{\unsure{exigences}} \inserted{propriétés} souhaitables qui seraient conséquences
\scanpage{72}\mspage{63}

de l'existence d'une telle fonction support, peuvent être \unsure{posés|posées} en \unsure{axiomes|axiome} sans inconvénient pratique.

La seule exception à cette règle concerne les géométries de lieux définis \deleted{\unsure{p}} en termes d'une fonction distance, où la disjonction est décrite en termes d'une distance $>\varepsilon$ \inserted{ou $\geq\varepsilon$}, et ne peut sans doute se ramener à une question de \deleted{\unsure{int|l'int}} ``disjonction ensembliste''. De toutes façons, je n'ai pas regardé ce \unsure{cas} d'assez près pour être assuré que le langage géométrique des magasins que je propose ici, s'y applique également.

\uline{Exemples} 1) \deleted{Dans} Dire que dans \unsure{les|des|le|un} magasin $M$, \deleted{\unsure{les}} $\lvert\unsure{\circ|0}\rvert$ est \unsure{définissable|definissable} par $\overset{\circ}{\ll}$, (p.~30, \unsure{31|71}) signifie : que
\[
 X \longrightarrow \mathrm{Omb}(X) \qquad M \longrightarrow \unsure{\mathcal{P}|P}(M)
\]
est une fonction support \inserted{\uline{fidèle}$^{*}$}. \deleted{\unsure{Mais cette|Mais alors|Mais elle}} \unsure{Sa|Ça} en fait un bon paquet !

% slanted
\margin{$^{*}$\unsure{NB} C'est \uuline{\unsure{tjrs}} une fonction support.}

% archive sheet number: 72
\scanpage{73}\mspage{64}

\uline{Exemple 2} Les magasins formés avec des \unsure{pseudosimplexes|pseudo-simplexes} \unsure{etc|etc.} dans des espaces \inserted{$X$} top. ou modérés (disons) --- \unsure{comme} ceux associés à une ``corde'' (ou ``ficelle'') topologique, ou plus gén., à un \unsure{ens|ens.} ordonné quelconque, échappent au cas précédent --- des lieux différents peuvent ne pas être disjoints. Mais pourtant, la fonction support fidèle à valeurs dans $\unsure{\mathcal{P}|P}(X)$ est \unsure{toute} trouvée !

\uline{Exemple 3} Magasins ``triviaux'' \unsure{ou|---} \unsure{maquettes|coquettes|casquettes}. La \deleted{condition} \inserted{relation} $X\lvert\unsure{\circ|0}\rvert Y$ signifie $X\ne Y$, la \deleted{condition} \inserted{relation} $X\overset{\circ}{\ll}Y$ signifie $X=Y$, on voit \deleted{que} donc que $X\lvert\unsure{\bullet|\circ}\rvert Y$ \unsure{sss|ssi} $\nexists Z$, $Z\overset{\circ}{\ll}X$, $Z\overset{\circ}{\ll}Y$, i.e. on est dans le cas de l'exemple 1. En associant à $\unsure{\text{tt}|\forall}\,X$
\[
 \lvert X\rvert\overset{\unsure{\mathrm{def}|\text{déf}}}{=}\widetilde{X}=\unsure{\dot M|M}_{\leq X}=\mathrm{Omb}(X)
\]
(ou encore, l'adhérence de $X$ pour la top. de $M$ associée à $\leq$), on trouve une fonction support fidèle. Notons qu'ici
\[
 \lvert X\rvert^{\circ}=\{X\} \qquad \text{pour tt}\ X\in M
\]

% archive sheet number: 73
